
\documentclass[11pt,a4paper]{article}

\usepackage[margin=1.7cm]{geometry}
\usepackage[T1]{fontenc}
\usepackage[utf8]{inputenc}
\usepackage{lmodern}
\usepackage{amsmath,amssymb}
\usepackage{booktabs}
\usepackage{array}
\usepackage{xcolor}
\usepackage{enumitem}
\usepackage{listings}
\usepackage{fancyhdr}
\usepackage{tabularx}

\definecolor{ForestGreen}{RGB}{22,100,70}
\definecolor{SoftGreen}{RGB}{235,246,239}
\definecolor{CodeGray}{RGB}{248,249,248}

\pagestyle{fancy}
\fancyhf{}
\lhead{\small Data Structures and Algorithms}
\rhead{\small Week 3 Practice Exercises}
\cfoot{\small \thepage}
\renewcommand{\headrulewidth}{0.4pt}

\setlength{\parindent}{0pt}
\setlength{\parskip}{4pt}
\setlist[enumerate]{leftmargin=1.6em,itemsep=2pt,topsep=3pt}
\setlist[itemize]{leftmargin=1.5em,itemsep=2pt,topsep=3pt}

\lstdefinestyle{pseudo}{
  basicstyle=\ttfamily\small,
  keywordstyle=\color{blue!70!black}\bfseries,
  morekeywords={FUNCTION,INPUT,OUTPUT,FOR,EACH,TO,DO,IF,THEN,ELSE,
    WHILE,RETURN,DIV,AND,OR},
  numbers=none,
  showstringspaces=false,
  keepspaces=true,
  columns=fullflexible,
  breaklines=true,
  frame=single,
  rulecolor=\color{ForestGreen!55},
  backgroundcolor=\color{CodeGray},
  tabsize=4
}

\newcommand{\exercise}[1]{%
  \vspace{0.7em}
  {\large\bfseries\color{ForestGreen}#1}
  \vspace{0.2em}
  \hrule
  \vspace{0.5em}
}

\begin{document}

\begin{center}
    {\LARGE\bfseries\color{ForestGreen} WEEK 3 PRACTICE EXERCISES}\\[0.3em]
    {\large Algorithms, Pseudocode, and Counting Basic Operations}\\[0.4em]
    \textbf{Sections 1--3: Definition of Algorithms, Algorithm Representation, and Complexity of Algorithms}
\end{center}

\vspace{0.4em}

\begin{center}
\begin{tabularx}{0.96\textwidth}{>{\bfseries}l X >{\bfseries}l X}
Name: & \hrulefill & Student ID: & \hrulefill
\end{tabularx}
\end{center}

\vspace{0.5em}

\colorbox{SoftGreen}{
\parbox{0.94\textwidth}{
\textbf{Counting convention.}
Unless stated otherwise, count each comparison, assignment, or execution of a statement as one basic operation.
For \texttt{FOR} loops, do not count loop-control operations separately.
The main goal is to determine how many times important statements are executed and then build the running-time function $T(n)$.
}}

\exercise{Exercise 1. Single loop with a condition}

\begin{lstlisting}[style=pseudo]
FUNCTION COUNT_POSITIVE
INPUT: array A of length n
OUTPUT: number of positive values

    count = 0

    FOR i = 0 TO n - 1 DO
        IF A[i] > 0 THEN
            count = count + 1

    RETURN count
\end{lstlisting}

\begin{enumerate}
    \item How many times is the condition $A[i] > 0$ evaluated?

    \item Suppose exactly $p$ elements of $A$ are positive.
    How many times is \texttt{count = count + 1} executed?

    \item Count the following basic operations:
    \begin{itemize}
        \item \texttt{count = 0};
        \item comparison \texttt{A[i] > 0};
        \item \texttt{count = count + 1};
        \item \texttt{RETURN count}.
    \end{itemize}
    Write the running-time function $T(n,p)$.

    \item Find $T_{\mathrm{best}}(n)$ and $T_{\mathrm{worst}}(n)$.
    Clearly state what kind of input gives each case.
\end{enumerate}

\exercise{Exercise 2. Find the maximum: same input size, different inputs}

\begin{lstlisting}[style=pseudo]
FUNCTION FIND_MAX
INPUT: non-empty array A of length n
OUTPUT: largest value

    maximum = A[0]

    FOR i = 1 TO n - 1 DO
        IF A[i] > maximum THEN
            maximum = A[i]

    RETURN maximum
\end{lstlisting}

\begin{enumerate}
    \item How many times is the comparison $A[i] > maximum$ evaluated?

    \item Let $u$ be the number of times the assignment
    \texttt{maximum = A[i]} is executed. Write $T(n,u)$.

    \item Calculate the number of basic operations for
    \[
        A=[10,8,7,3,1].
    \]

    \item Calculate the number of basic operations for
    \[
        A=[1,3,5,8,10].
    \]

    \item Which type of input gives the worst case?
    Write $T_{\mathrm{worst}}(n)$.
\end{enumerate}

\exercise{Exercise 3. Two nested loops}

\begin{lstlisting}[style=pseudo]
FUNCTION COUNT_PAIRS
INPUT: array A of length n

    count = 0

    FOR i = 0 TO n - 1 DO
        FOR j = 0 TO n - 1 DO
            count = count + 1

    RETURN count
\end{lstlisting}

\begin{enumerate}
    \item For a fixed value of \texttt{i}, how many times does the inner loop execute?

    \item How many different values does \texttt{i} take?

    \item Complete the following table.

    \begin{center}
    \renewcommand{\arraystretch}{1.35}
    \begin{tabular}{c|c|c|c}
        \toprule
        $n$ & Inner loop per $i$ & Number of values of $i$ & Total executions \\
        \midrule
        $2$ & & & \\
        $3$ & & & \\
        $5$ & & & \\
        $n$ & & & \\
        \bottomrule
    \end{tabular}
    \end{center}

    \item How many times is \texttt{count = count + 1} executed?

    \item Counting initialization and return as one operation each, find $T(n)$.
\end{enumerate}

\exercise{Exercise 4. Nested loops with changing bounds}

\begin{lstlisting}[style=pseudo]
FUNCTION COUNT_UNORDERED_PAIRS
INPUT: array A of length n

    count = 0

    FOR i = 0 TO n - 2 DO
        FOR j = i + 1 TO n - 1 DO
            count = count + 1

    RETURN count
\end{lstlisting}

\begin{enumerate}
    \item Complete the table below.

    \begin{center}
    \renewcommand{\arraystretch}{1.35}
    \begin{tabular}{c|c}
        \toprule
        Value of $i$ & Number of values of $j$ \\
        \midrule
        $0$ & \\
        $1$ & \\
        $2$ & \\
        $\vdots$ & $\vdots$ \\
        $n-2$ & \\
        \bottomrule
    \end{tabular}
    \end{center}

    \item Show that the total number of executions of
    \texttt{count = count + 1} is
    \[
        (n-1)+(n-2)+\cdots+2+1.
    \]

    \item Use
    \[
        1+2+\cdots+(n-1)=\frac{n(n-1)}{2}
    \]
    to obtain the total number of executions.

    \item Counting initialization and return as one operation each, find $T(n)$.

    \item Verify your result manually for $n=4$ by listing all pairs $(i,j)$ processed by the algorithm.
\end{enumerate}

\exercise{Exercise 5. Repeated halving}

\begin{lstlisting}[style=pseudo]
FUNCTION HALVING
INPUT: positive integer n

    remaining = n
    count = 0

    WHILE remaining > 1 DO
        remaining = remaining DIV 2
        count = count + 1

    RETURN count
\end{lstlisting}

Assume first that
\[
    n=2^k.
\]

\begin{enumerate}
    \item Trace the algorithm for $n=16$.

    \begin{center}
    \renewcommand{\arraystretch}{1.35}
    \begin{tabular}{c|c|c}
        \toprule
        Iteration & \texttt{remaining} before & \texttt{remaining} after \\
        \midrule
        $1$ & & \\
        $2$ & & \\
        $3$ & & \\
        $4$ & & \\
        \bottomrule
    \end{tabular}
    \end{center}

    \item How many loop iterations are performed for:
    \[
        n=8,\qquad16,\qquad32,\qquad64?
    \]

    \item Complete:
    \[
        n=2^k
        \qquad\Longrightarrow\qquad
        k=\underline{\hspace{2.5cm}}.
    \]

    \item Hence, how many times does the loop body execute as a function of $n$?

    \item Now count:
    \begin{itemize}
        \item \texttt{remaining = n}: 1 operation;
        \item \texttt{count = 0}: 1 operation;
        \item each test of \texttt{remaining > 1}: 1 operation;
        \item each statement inside the loop body: 1 operation;
        \item \texttt{RETURN count}: 1 operation.
    \end{itemize}

    Find the exact running-time function $T(n)$ for $n=2^k$.
\end{enumerate}

\vfill

\begin{center}
\color{ForestGreen}\bfseries
Focus on the counting process:
Algorithm $\rightarrow$ number of executions $\rightarrow$ $T(n)$.
\end{center}

\end{document}
